% Candidate Section 2 text accompanying figures.tex.
\section{Action Variability in Robot Demonstrations}
\label{sec:data-structure}

\subsection{Action scale changes across task phases}

Robot demonstrations do not share a single action scale. We first use a
converged MSE policy as an operational estimate of the conditional mean and
measure its residual RMS separately for every input. The resulting scale
distributions are broad (Figure~\ref{fig:data-phase-scale}, left): their
10th--90th-percentile spans range from $1.8\times$ to $4.6\times$, and the full
GR1/RoboCasa range reaches $51\times$. Because these residuals may include model
error, we independently test the same property using demonstrations alone. For
each task stage, we estimate an action radius from one fold of demonstrations
and evaluate its residual variance on a disjoint fold. Across 149 tasks from
five datasets, high-radius stages have $1.8$--$4.6\times$ the held-out variance
of low-radius stages (Figure~\ref{fig:data-phase-scale}, right). A fixed-variance
Gaussian assigns the same scale to all of these inputs; heteroscedastic
regression instead learns the required observation-dependent $\sigma(o)$.

\subsection{Local residuals are heavy-tailed}

Input-dependent scale does not by itself capture the shape of the residual
distribution. To isolate shape from scale, we standardize each residual by its
local RMS and additionally normalize every action coordinate using statistics
estimated only from the training demonstrations. The held-out residuals still
depart systematically from a Gaussian (Figure~\ref{fig:data-heavy-tail}):
events beyond three standard deviations occur 3.6 times as often as predicted
by the Gaussian fit. A fitted Student-$t$ distribution follows this tail and
achieves lower held-out negative log-likelihood. These two observations
motivate a heteroscedastic Student-$t$ objective: the predicted scale adapts
across task stages, while the heavy-tailed likelihood limits the influence of
rare deviations.
